\documentclass[12pt, a4paper]{article}


\usepackage{xcolor}
\usepackage{geometry}
\usepackage{hyperref}
\usepackage{titlesec}
\usepackage{enumitem}
\usepackage{longtable}
\usepackage{booktabs}
\usepackage{array}
\usepackage{ragged2e}
\usepackage{fancyhdr}
% ---------- پکیج‌های پایه ----------
\usepackage{xepersian}
\settextfont{Vazirmatn}
\setlatintextfont{Times New Roman}
\setdigitfont{Vazirmatn}

\usepackage{geometry}
\geometry{margin=2.5cm}

\usepackage{hyperref}
\hypersetup{
    colorlinks=true,
    linkcolor=blue,
    urlcolor=blue,
    citecolor=blue
}

\usepackage{titlesec}
\titleformat{\section}{\Large\bfseries\color{blue!60!black}}{\thesection}{1em}{}
\titleformat{\subsection}{\large\bfseries}{\thesubsection}{1em}{}
\titleformat{\subsubsection}{\normalsize\bfseries}{\thesubsubsection}{1em}{}

\usepackage{enumitem}
\setlist{nosep, leftmargin=*}

\usepackage{longtable}
\usepackage{booktabs}
\usepackage{array}
\usepackage{ragged2e}

\usepackage{fancyhdr}
\pagestyle{fancy}
\fancyhf{}
\fancyhead[R]{\small سند چشم‌انداز — هوم‌استور}
\fancyhead[L]{\small نسخهٔ ۱.۰}
\fancyfoot[C]{\thepage}
\renewcommand{\headrulewidth}{0.4pt}

\usepackage{xcolor}
\definecolor{primary}{RGB}{139, 90, 43}
\definecolor{accent}{RGB}{224, 122, 63}

\linespread{1.4}

% ---------- اطلاعات سند ----------
\title{
    \vspace{-1cm}
    \textbf{\Huge سند چشم‌انداز}\\[0.3cm]
    \textbf{\LARGE پروژهٔ هوم‌استور}\\[0.5cm]
    \large یک فروشگاه آنلاین لوازم خانگی
}
\author{
    تهیه‌کننده: فضلی\\
    دانشگاه: أزاد واحد زنجان\\
    ترم: اول ۱۴۰۵
}
\date{\today}

% ============================================================
\begin{document}

\maketitle
\thispagestyle{empty}

\vspace{1cm}

\begin{center}
\begin{tabular}{|l|l|}
\hline
\textbf{نام پروژه} & هوم‌استور (HomeStore) \\
\hline
\textbf{نسخهٔ سند} & ۱.۰ \\
\hline
\textbf{تاریخ} & \today \\
\hline
\textbf{تهیه‌کننده} & فضلی \\
\hline
\textbf{وضعیت} & پیش‌نویس \\
\hline
\end{tabular}
\end{center}

\newpage
\tableofcontents
\newpage

% ============================================================
\section{مقدمه}

\subsection{هدف سند}
این سند چشم‌انداز، دامنه و ذی‌نفعان پروژهٔ \textbf{هوم‌استور} را تعریف می‌کند.
هدف آن ایجاد یک درک مشترک میان تیم توسعه، مشتری و ذی‌نفعان درباره
«چه چیزی می‌سازیم، برای چه کسی، و چرا» است.

این سند \textbf{قرارداد رسمی} نیست، بلکه \textbf{مرجع راهنما} برای تصمیم‌گیری‌های
بعدی در طول چرخهٔ حیات پروژه است و در هر Sprint بازبینی می‌شود.

\subsection{دامنهٔ سند}
\begin{itemize}
    \item \textbf{در دامنه:} اهداف کسب‌وکار، کاربران هدف، محدودهٔ MVP، معیارهای موفقیت، ریسک‌های سطح بالا.
    \item \textbf{خارج از دامنه:} جزئیات فنی، طراحی UI، نیازمندی‌های دقیق functional.
\end{itemize}

\subsection{تعاریف و اصطلاحات}
\begin{center}
\begin{tabular}{|c|p{9cm}|}
\hline
\textbf{اصطلاح} & \textbf{تعریف} \\
\hline
MVP & Minimum Viable Product — کمینهٔ محصول قابل ارائه \\
\hline
User Story & توصیف نیازمندی از دید کاربر \\
\hline
Persona & نمایندهٔ فرضی کاربر هدف \\
\hline
Sprint & بازهٔ کوتاه توسعه (۲-۳ هفته) در Scrum \\
\hline
DoD & Definition of Done — معیار «تمام‌شده» بودن کار \\
\hline
\end{tabular}
\end{center}

\subsection{مراجع}
\begin{itemize}
    \item IEEE 830 — Recommended Practice for Software Requirements Specifications
    \item Roman Pichler — \textit{Product Vision Board}
    \item Scrum Guide (2020)
    \item Sommerville — \textit{Software Engineering}, Chapters 1-2
\end{itemize}

% ============================================================
\section{چشم‌انداز محصول}

\subsection{بیانیهٔ چشم‌انداز}

\begin{quote}
\textbf{برای} صاحبان فروشگاه‌های کوچک لوازم خانگی \textbf{که} محصولات اصل و
با گارانتی می‌فروشند اما فقط مشتریان محلی آن‌ها را می‌شناسند،
\textbf{هوم‌استور} یک \textbf{فروشگاه آنلاین} است
\textbf{که} کاتالوگ، مقایسه و ثبت سفارش را به‌صورت ساده و قابل‌اعتماد فراهم می‌کند.

\textbf{برخلاف} مارکت‌پلیس‌های بزرگ که فروشنده‌های متعدد و کیفیت نامطمئن دارند،
\textbf{محصول ما} روی یک فروشندهٔ قابل‌اعتماد با گارانتی واقعی تمرکز دارد.
\end{quote}

\subsection{ارزش پیشنهادی}
\begin{center}
\begin{tabular}{|l|p{8cm}|}
\hline
\textbf{برای} & \textbf{ارزش} \\
\hline
مشتری نهایی & خرید سریع، مقایسهٔ آسان، اطمینان از گارانتی \\
\hline
صاحب فروشگاه (ادمین) & دسترسی به مشتریان فراتر از محله، مدیریت سادهٔ محصولات \\
\hline
انباردار & ثبت و پیگیری دقیق موجودی بدون کاغذبازی \\
\hline
\end{tabular}
\end{center}

\subsection{مزیت رقابتی}
\begin{itemize}
    \item \textbf{سادگی:} رابط کاربری ساده‌تر از دیجی‌کالا و بامیلو
    \item \textbf{تمرکز:} فقط لوازم خانگی، نه هر چیز
    \item \textbf{اعتماد:} گارانتی و اصالت محصول تضمین‌شده
    \item \textbf{محلی‌بودن:} پشتیبانی از مشتریان شهرهای کوچک
\end{itemize}

% ============================================================
\section{مشکل و فرصت}

\subsection{مشکل فعلی}
صاحبان فروشگاه‌های کوچک لوازم خانگی با چالش‌های زیر روبه‌رو هستند:
\begin{enumerate}
    \item \textbf{محدودیت جغرافیایی:} فقط مشتریان محلی از آن‌ها خرید می‌کنند.
    \item \textbf{نبود کانال آنلاین:} نمی‌توانند محصولات خود را آنلاین عرضه کنند.
    \item \textbf{فرآیند دستی سفارش:} سفارش‌ها تلفنی، یادداشت‌شده و پرخطا.
    \item \textbf{عدم شفافیت موجودی:} مشتری نمی‌داند چه چیزی موجود است.
    \item \textbf{نبود امکان مقایسه:} مشتری نمی‌تواند محصولات را کنار هم ببیند.
\end{enumerate}

\subsection{فرصت}
\begin{itemize}
    \item رشد خرید آنلاین لوازم خانگی در ایران
    \item افزایش نفوذ اینترنت در شهرهای کوچک
    \item عدم حضور فروشگاه‌های کوچک در فضای آنلاین
    \item امکان ایجاد یک بازار نیچ برای محصولات اصل با گارانتی
\end{itemize}

% ============================================================
\section{کاربران هدف و ذی‌نفعان}

\subsection{ذی‌نفعان}
\begin{center}
\begin{tabular}{|l|l|p{5cm}|}
\hline
\textbf{ذی‌نفع} & \textbf{نقش} & \textbf{علاقه} \\
\hline
خانم رضایی & صاحب فروشگاه (ادمین) & فروش بیشتر، مدیریت آسان \\
\hline
مشتری نهایی & خریدار لوازم خانگی & خرید سریع، اطمینان، مقایسه \\
\hline
انباردار & مسئول موجودی & ثبت دقیق، دسترسی سریع \\
\hline
تیم توسعه & سازندهٔ محصول & تحویل به‌موقع، یادگیری \\
\hline
استاد درس & ناظر پروژه & رعایت اصول مهندسی \\
\hline
\end{tabular}
\end{center}

\subsection{پرسوناهای کلیدی}
جزئیات کامل در \texttt{personas.md}. خلاصه:
\begin{center}
\begin{tabular}{|l|l|l|}
\hline
\textbf{پرسونا} & \textbf{نقش} & \textbf{ویژگی کلیدی} \\
\hline
خانم رضایی & ادمین & ۴۲ ساله، مهارت فنی متوسط \\
\hline
مریم & مشتری & ۳۲ ساله، خریدار آگاه \\
\hline
رضا & انباردار & ۲۸ ساله، مهارت فنی پایین \\
\hline
\end{tabular}
\end{center}

% ============================================================
\section{محدودهٔ محصول}

\subsection{داخل محدوده (MVP)}

\textbf{کاربری (Functional):}
\begin{itemize}
    \item \textbf{F-01} نمایش کاتالوگ محصولات با تصویر، قیمت و مشخصات
    \item \textbf{F-02} جست‌وجو و فیلتر پیشرفته (برند، قیمت، مصرف انرژی، ظرفیت)
    \item \textbf{F-03} صفحهٔ جزئیات محصول با گارانتی و مشخصات فنی
    \item \textbf{F-04} مقایسهٔ ۲-۳ محصول در جدول
    \item \textbf{F-05} سبد خرید (افزودن، حذف، تغییر تعداد)
    \item \textbf{F-06} ثبت سفارش با فرم و اعتبارسنجی
    \item \textbf{F-07} پنل ادمین: افزودن/ویرایش/حذف محصول
    \item \textbf{F-08} پنل ادمین: مشاهدهٔ سفارش‌ها و تغییر وضعیت
\end{itemize}

\textbf{غیرکاربری (Non-functional):}
\begin{itemize}
    \item \textbf{NF-01} زمان بارگذاری صفحه < ۲ ثانیه
    \item \textbf{NF-02} پشتیبانی از ۱۰۰ کاربر همزمان
    \item \textbf{NF-03} رابط کاربری واکنش‌گرا برای موبایل و دسکتاپ
    \item \textbf{NF-04} رمزنگاری گذرواژه‌ها (bcrypt)
    \item \textbf{NF-05} پشتیبانی از زبان فارسی و RTL
\end{itemize}

\subsection{خارج از محدوده}
\begin{center}
\begin{tabular}{|l|p{8cm}|}
\hline
\textbf{مورد} & \textbf{دلیل} \\
\hline
پرداخت واقعی (درگاه بانکی) & نیاز به مجوز و زیرساخت — فقط شبیه‌سازی \\
\hline
اپلیکیشن موبایل & تمرکز روی وب \\
\hline
سیستم حمل و نقل & خارج از توان یک ترم \\
\hline
گارانتی واقعی و خدمات پس از فروش & پیچیدگی حقوقی \\
\hline
چت پشتیبانی آنلاین & اولویت پایین \\
\hline
چند زبانه & تمرکز روی فارسی \\
\hline
\end{tabular}
\end{center}

\subsection{محدودیت‌ها}
\begin{itemize}
    \item \textbf{زمانی:} ۱۴ هفته، ۳ Sprint
    \item \textbf{تیمی:} ۳-۴ نفر
    \item \textbf{فنی:} بدون سرور اختصاصی، استقرار روی سرویس ابری رایگان
    \item \textbf{بودجه:} صفر (فقط ابزارهای رایگان)
\end{itemize}

% ============================================================
\section{معیارهای موفقیت}

\subsection{معیارهای کاربری}
\begin{center}
\begin{tabular}{|l|l|l|}
\hline
\textbf{معیار} & \textbf{هدف} & \textbf{نحوهٔ اندازه‌گیری} \\
\hline
زمان ثبت سفارش & < ۳ دقیقه & تست کاربری \\
\hline
زمان افزودن محصول & < ۲ دقیقه & تست کاربری \\
\hline
رضایت کاربر & > ۸۰٪ & نظرسنجی ساده \\
\hline
\end{tabular}
\end{center}

\subsection{معیارهای فنی}
\begin{center}
\begin{tabular}{|l|l|}
\hline
\textbf{معیار} & \textbf{هدف} \\
\hline
پوشش تست (Code Coverage) & > ۷۰٪ \\
\hline
تعداد باگ بحرانی در انتشار & صفر \\
\hline
زمان Build در CI & < ۵ دقیقه \\
\hline
موفقیت Deploy & ۱۰۰٪ \\
\hline
\end{tabular}
\end{center}

\subsection{معیارهای فرآیندی}
\begin{itemize}
    \item تحویل هر Sprint به‌موقع
    \item رعایت Definition of Done در هر Story
    \item مستندات به‌روز در پایان هر Sprint
\end{itemize}

% ============================================================
\section{ریسک‌های سطح بالا}

\begin{center}
\begin{tabular}{|p{4cm}|c|c|p{4cm}|}
\hline
\textbf{ریسک} & \textbf{احتمال} & \textbf{اثر} & \textbf{راهکار کاهش} \\
\hline
تغییر نیازمندی‌ها & بالا & متوسط & Sprintهای کوتاه، Backlog منعطف \\
\hline
کمبود زمان & متوسط & بالا & MVP کوچک، اولویت MoSCoW \\
\hline
ضعف فنی تیم & متوسط & متوسط & منتورینگ استاد، تقسیم کار \\
\hline
از دست دادن داده & کم & بالا & Backup خودکار، Git \\
\hline
Scope Creep & بالا & بالا & مرزبندی دقیق MVP، تأیید استاد \\
\hline
\end{tabular}
\end{center}

% ============================================================
\section{رویکرد و متدولوژی}

\subsection{متدولوژی}
\begin{itemize}
    \item \textbf{Scrum سبک‌شده:} ۳ Sprint دو-سه هفته‌ای
    \item \textbf{ابزار مدیریت:} GitHub Projects
    \item \textbf{جلسات:} Sprint Planning، Daily کوتاه، Review، Retrospective
\end{itemize}

\subsection{فناوری پیشنهادی}
\begin{center}
\begin{tabular}{|l|l|}
\hline
\textbf{لایه} & \textbf{فناوری} \\
\hline
Backend & Node.js + Express \\
\hline
Frontend & React \\
\hline
Database & PostgreSQL \\
\hline
Version Control & Git + GitHub \\
\hline
CI/CD & GitHub Actions \\
\hline
مستندسازی & Markdown در GitHub \\
\hline
\end{tabular}
\end{center}

\subsection{Definition of Done}
یک Story وقتی «تمام» است که:
\begin{enumerate}
    \item کد نوشته و Merge شده باشد.
    \item تست Unit برای آن نوشته شده باشد.
    \item Code Review شده باشد.
    \item مستندات مربوطه به‌روز شده باشد.
    \item در محیط Demo کار کند.
\end{enumerate}

% ============================================================
\section{نقشهٔ راه}

\begin{center}
\begin{tabular}{|c|c|p{4cm}|p{4.5cm}|}
\hline
\textbf{Sprint} & \textbf{بازه} & \textbf{تمرکز} & \textbf{خروجی} \\
\hline
Sprint 1 & هفته ۱-۵ & کاتالوگ و کشف & لیست محصولات، جست‌وجو، فیلتر، سبد خرید \\
\hline
Sprint 2 & هفته ۶-۱۱ & تراکنش و مدیریت & ثبت سفارش، مقایسه، پنل ادمین \\
\hline
Sprint 3 & هفته ۱۲-۱۴ & کیفیت و انتشار & تست کامل، CI/CD، مستندات، Demo Day \\
\hline
\end{tabular}
\end{center}

% ============================================================
\section{تأیید و تصویب}

\begin{center}
\begin{tabular}{|l|l|l|l|}
\hline
\textbf{نقش} & \textbf{نام} & \textbf{تاریخ} & \textbf{امضا} \\
\hline
صاحب محصول & [نام استاد] & \today & \\
\hline
مشتری فرضی & خانم رضایی & \today & \\
\hline
نمایندهٔ تیم & [نام] & \today & \\
\hline
\end{tabular}
\end{center}

% ============================================================
\appendix
\section{تاریخچهٔ تغییرات}

\begin{center}
\begin{tabular}{|c|c|l|l|}
\hline
\textbf{نسخه} & \textbf{تاریخ} & \textbf{تغییرات} & \textbf{تهیه‌کننده} \\
\hline
۱.۰ & \today & نسخهٔ اولیه & [نام] \\
\hline
 & & & \\
\hline
\end{tabular}
\end{center}

\end{document}
